% GENERATED, edit the French unit or the translation table
\mxunit{premiere-c09}{The exponential function and continuous models}\label{C09}
\begin{keypoint}[Chapter path]
\textbf{Objectives.} Use the algebraic properties; differentiate; study growth or decay; link the exponential to geometric sequences.

\textbf{Prerequisites.} Powers, geometric sequences, differentiation and sign.

\textbf{Yearly reference point.} 2--11 February: 8 hours\end{keypoint}
\section*{Diagnostic: six questions, 10 minutes}\label{C09-D}
Answer without consulting the lesson. Mark the questions for which help is needed. This diagnostic is not graded.
\begin{enumerate}
\item Simplify $2^3\times2^4$.
\item Express $3^{-2}$.
\item Solve $2x+1=5$.
\item Differentiate $x^2+2x$.
\item A \pc{20} decrease gives which coefficient?
\item For a geometric sequence $u_0=5,q=2$, calculate $u_n$.
\end{enumerate}
\paragraph{After the diagnostic.} Revisit the flagged topics with the teacher; redo a similar question, then explain the rule used.
\paragraph{Speaking and logic.} Explain the difference between a continuous rate and a discrete multiplying coefficient; present an inequality proved through the minimum of a difference.
\mxnewpage
\section*{Investigate before formalising}\label{C09-RECH}
Compare the increments of $e^x$ and the values of the function around 0, with several positive steps. Build three Euler steps and explain the rule chosen.
\paragraph{Organisation.} Individual investigation 5 min, comparison in pairs 5 min, class discussion 5 min. Keep a record of one attempt, even an unfinished one. A conjecture must then be justified.
\mxnewpage
\section{Defining a new function}\label{C09-U1}
\begin{keypoint}[Definition and results]
The exponential function is the unique function f differentiable on $\R$ satisfying $f\prime=f$ and $f(0)=1$. Its existence and uniqueness are accepted. It is written $\exp$, then $e^x$ with $e=\exp(1)\approx\num{2,71828}$. The number e is a constant; x is the variable. The exponential is strictly positive.
\end{keypoint}
\paragraph{Status and justification.} Existence and uniqueness accepted in accordance with the programme. Positivity is deduced, after the algebraic properties accepted below, from $\exp (x)=\exp (\dfrac{x}{2})^2$ and from $\exp(x)\exp(-x)=1$.
\begin{example}$\exp(0)=1$, $\exp(1)=e$, $\exp\prime(0)=1$. The tangent at 0 is therefore $y=1+x$. The notation $e^x$ extends the properties of powers to every real exponent.\end{example}\label{C09-EX1}
\mxinclude{figures/C09-U1}
\begin{pitfall}$e^0=1$, not 0. The function never takes the value 0, even for negative x.\end{pitfall}
\begin{application}What is the equation of the tangent to $e^x$ at 1?\end{application}\label{C09-AP1}
\section{Algebraic properties and comparisons}\label{C09-U2}
\begin{keypoint}[Definition and results]
For all real x,y, $e^{x+y}=e^xe^y$, $e^{-x}=\dfrac{1}{e^x}$, $e^{x-y}=\dfrac{e^x}{e^y}$. For integer n, $(e^x)^n=e^{nx}$. The function is strictly increasing on $\R$, so $e^u=e^v\iff u=v$ and $e^u<e^v\iff u<v$.
\end{keypoint}
\paragraph{Status and justification.} The addition relation is accepted in the common pathway; it can be explored further via uniqueness. Positivity and $\exp\prime=\exp>0$ give strict increase. This justifies the equivalences for equations and inequalities.
\begin{example}$e^{2x}e^{1-x}=e^{x+1}$. $e^{3x-2}=e^{x+4}$ is equivalent to $3x-2=x+4$, so x=3. $e^{2x}\leq1=e^0$ is equivalent to $x\leq0$.\end{example}\label{C09-EX2}
\begin{pitfall}$e^{x+y}$ is a product, in general never $e^x+e^y$. We do not solve $e^x=2$ with a logarithm, which is not studied in Première.\end{pitfall}
\begin{application}Simplify $\dfrac{e^{3x}}{e^{x-1}}$; solve $e^{x+2}>e^{2x}$.\end{application}\label{C09-AP2}
\section{Differentiating and studying}\label{C09-U3}
\begin{keypoint}[Definition and results]
For real a, $t\mapsto e^{at}$ is differentiable and its derivative is $ae^{at}$. For $e^{at+b}=e^be^{at}$, the derivative is $ae^{at+b}$. The product and quotient rules apply with the derivatives already known. The sign of the exponential often simplifies the study of the sign of the derivative.
\end{keypoint}
\paragraph{Status and justification.} The formula for differentiating $e^{at}$ is accepted; the case with b follows by taking out the constant factor $e^b$. The general chain rule is not needed here.
\begin{example}For $f(x)=(x+1)e^{-x}$, $f\prime(x)=e^{-x}-(x+1)e^{-x}=-xe^{-x}$. Since $e^{-x}>0$, the derivative has the sign of $-x$: global maximum 1 at 0.\end{example}\label{C09-EX3}
\begin{pitfall}Do not forget the factor a or the negative sign in $e^{-x}$. The exponential factor is positive, not necessarily increasing.\end{pitfall}
\begin{application}Differentiate $g(t)=3e^{-2t}$ and study its direction of variation.\end{application}\label{C09-AP3}
\section{From discrete to continuous}\label{C09-U4}
\begin{keypoint}[Definition and results]
A model $Q(t)=Q_0e^{kt}$, with $Q_0>0$, represents growth if $k>0$, decay if $k<0$, constancy if $k=0$. For a fixed step h, $\dfrac{Q(t+h)}{Q(t)}=e^{kh}$. The values at integer times form a geometric sequence with common ratio $e^k$. The unit of k is the inverse of the unit of time.
\end{keypoint}
\paragraph{Status and justification.} The ratio follows from $e^{k(t+h)}=e^{kt}e^{kh}$. The direction of variation is deduced from $Q\prime(t)=kQ_0e^{kt}$ and the positivity of the remaining factor.
\begin{example}$Q(t)=100e^{-\num{0,2}t}$, t in hours: $Q(0)=100$, $\dfrac{Q(t+1)}{Q(t)}=e^{-\num{0,2}}$. A simulation can look for the first whole number of hours at which Q falls below 50; the exact equation is not solved with a logarithm in this chapter.\end{example}\label{C09-EX4}
\mxinclude{figures/C09-U4}
\begin{pitfall}The relative rate per discrete unit of time is $e^k-1$, not exactly k; they are close only for small k.\end{pitfall}
\begin{application}For $Q(t)=5e^{\num{0,3}t}$, express $Q(n)$ and $\dfrac{Q(t+2)}{Q(t)}$.\end{application}\label{C09-AP4}
\section{Building numerically and proving an inequality}\label{C09-U5}
\begin{keypoint}[Definition and results]
To approximate the solution of $y\prime=y$, $y(0)=1$ with a step $h$, Euler replaces the increment by $hy$: $y_{k+1}=(1+h)y_k$. On $\intff{0}{1}$ with $h=\dfrac1n$, the approximation of $e$ is $\left(1+\dfrac1n\right)^n$. This construction is experimental. The inequality $e^x\geq1+x$ is obtained by studying a function.
\end{keypoint}
\paragraph{Status and justification.} For $g(x)=e^x-1-x$, $g\prime(x)=e^x-1$, negative if x<0 and positive if x>0. The minimum is $g(0)=0$, which proves the inequality; equality only at 0. The convergence of Euler's method is not proved here.
\begin{example}With n=2, Euler gives $y_0=1,y_1=\num{1,5},y_2=\num{2,25}$. With n=10, $(\num{1,1})^{10}\approx\num{2,59374}$. We compare with e; we do not confuse an approximate value with the definition of the function.\end{example}\label{C09-EX5}
\begin{pitfall}A graph showing a curve above a line does not prove the inequality for all real numbers.\end{pitfall}
\begin{application}Show that $e^{2x}\geq1+2x$ for every real number x.\end{application}\label{C09-AP5}
\mxnewpage
\section*{Exercise path}\label{C09-PARC}
Start with exercises 1 and 2. Then choose one of exercises 3 or 4 and do it without help. Exercises 5 to 7 are a reserve for transfer, investigation or programming: the teacher designates one. Exercise 8 is revisited after a delay.
\begingroup
\mxexercisehead{C09-E1}{Transforming}
For every real number $x$, consider the expressions\[A=e^xe^{2-x},\qquad B=\dfrac{(e^{2x})^3}{e^x},\qquad C=e^{-x}(e^x+1).\]
\begin{enumerate}[label=\arabic*.]
\item Simplify $A$ and $B$, justifying each grouping of exponents.
\item Expand then simplify $C$. Explain why a sum does not simplify in the same way as a product.
\end{enumerate}
\mxexercisehead{C09-E2}{Equations and inequalities}
Consider an equation and an inequality in the real unknown $x$:\[e^{2x-1}=e^{x+3}\qquad\text{and}\qquad e^{-3x}\geq e^{x-2}.\]
\begin{enumerate}[label=\arabic*.]
\item Solve each in $\R$, citing the property of the exponential that allows the exponents to be compared.
\item While solving the inequality, explain at which step the direction of the inequality changes. Give the solution set as an interval.
\end{enumerate}
\mxexercisehead{C09-E3}{Product and maximum}
The function $f$ is defined on $\intff{0}{4}$ by $f(x)=xe^{-x}$.
\begin{enumerate}[label=\arabic*.]
\item Calculate $f\prime(x)$, then study its sign on the interval.
\item Draw the variation table of $f$. Deduce its maximum and the x-coordinate where it is reached.
\item Check that the values at the endpoints do not give a larger maximum. Keep the exact values.
\end{enumerate}
\mxexercisehead{C09-E4}{A modelled population}
The following two parts are independent.\par\textbf{Part A: A population.} A population is modelled by $P(t)=200e^{\num{0,1}t}$, where $t\geq0$ is the time elapsed in years. $P(t)$ is a theoretical count, which is not necessarily an integer.
\begin{enumerate}[label=\arabic*.]
\item Determine the initial count, then the multiplying coefficient between times $t$ and $t+1$. Deduce the exact annual rate and its percentage rounded to the nearest hundredth.
\item Explain why this model cannot necessarily describe a real population indefinitely.
\item \textbf{Part B: Two curves.} In the same frame, sketch the curves of $t\mapsto e^{\dfrac{t}{2}}$ and $t\mapsto e^{-\dfrac{t}{2}}$. State their variation, their common point and their symmetry. Identify each curve clearly.
\end{enumerate}
\mxexercisehead{C09-E5}{Searching for a threshold}
The quantity of a substance, in milligrams, is modelled at time $t\geq0$ by $Q(t)=80e^{-\dfrac{t}{2}}$, where $t$ is expressed in hours. Checks take place at integer times $n=0,1,2,\ldots$. We look for the first check at which the quantity is strictly less than 10 mg.
\begin{enumerate}[label=\arabic*.]
\item Write an algorithm that searches for this first integer. Provide an iteration bound and state what the program returns if the bound is reached.
\item Run the algorithm, then check numerically the quantity at the index found and at the previous index.
\item Justify, using the decrease of $Q$, that this double check establishes that it is indeed the first index. Do not use a logarithm.
\end{enumerate}
\mxexercisehead{C09-E6}{A symmetry}
Consider the function $f$ defined on $\R$ by $f(x)=e^x+e^{-x}$.
\begin{enumerate}[label=\arabic*.]
\item Study the parity of $f$ and interpret the result for its curve.
\item Prove that $f(x)\geq2$ for every real number $x$, choosing a method. State the equality case.
\item Suggest a check of the result by a second method, from a square or a study of variation.
\end{enumerate}
\mxexercisehead{C09-E7}{Uniqueness under a hypothesis}
\textbf{Extension.} Assume that the exponential function has already been constructed. Let $f$ be a function differentiable on $\R$ that satisfies $f\prime(x)=f(x)$ for every real number $x$ and $f(0)=1$. Let $g(x)=f(x)e^{-x}$.
\begin{enumerate}[label=\arabic*.]
\item Calculate $g\prime$ and explain why $g$ is constant on $\R$.
\item Determine this constant, then deduce the expression of $f$.
\item Does this proof establish the existence or the uniqueness of the function sought? Justify.
\end{enumerate}
\mxexercisehead{C09-E8}{Delayed review}
Define the sequence $(u_n)$ by $u_n=e^{2n+1}$ for every natural number $n$.
\begin{enumerate}[label=\arabic*.]
\item Prove that this sequence is geometric, then state its first term and common ratio.
\item Calculate $u_0+u_1+u_2$ in exact form. Compare the direct calculation with the use of a sum formula.
\end{enumerate}
\paragraph{Homework.} Do only the task designated at the end of the session, in the dated progression: 10--20 min at most. The exercises depend on the topics already studied, so some questions are deferred. Revisit exercise 8 one week after the chapter (10 min). If the allotted time is exceeded, note the difficulty and bring your attempts; do not pile up the whole reserve.
\endgroup
\mxneed{12}\section*{To do in the CQFD textbook}
These exercises are done in your textbook. The progression indicates the session; they replace the homework planned for that day (10 to 20 minutes in total).
\paragraph{CQFD, exercise 18 p.~187, questions a, b and c.}\label{C09-M1} Simplify the products and quotients, citing the property of exponentials used.
\paragraph{CQFD, exercise 30 p.~187, questions a, b and c.}\label{C09-M2} Solve the three inequalities and write the solution sets as intervals.
\begingroup\small
\mxnewpage\section*{MCQ 1: understand and choose}\label{C09-Q1}
Five single-answer questions for practice. Choose one answer per question; then explain one distractor. This chapter format does not reproduce the official EAM format.
\begin{enumerate}
\item $e^0=$
\begin{enumerate}[label=\Alph*.,itemsep=2pt]
\item $0$
\item $1$
\item $e$
\item $-1$
\end{enumerate}
\item $e^xe^y=$
\begin{enumerate}[label=\Alph*.,itemsep=2pt]
\item $e^{x+y}$
\item $e^{xy}$
\item $e^x+e^y$
\item $e^{x-y}$
\end{enumerate}
\item Derivative of $e^{-2x}$:
\begin{enumerate}[label=\Alph*.,itemsep=2pt]
\item $e^{-2x}$
\item $2e^{-2x}$
\item $-e^{-2x}$
\item $-2e^{-2x}$
\end{enumerate}
\item $e^x=1$ is equivalent to
\begin{enumerate}[label=\Alph*.,itemsep=2pt]
\item $x=1$
\item $x=e$
\item $x=0$
\item no solution
\end{enumerate}
\item The sign of $e^{-x}$ is
\begin{enumerate}[label=\Alph*.,itemsep=2pt]
\item negative if x>0
\item always strictly positive
\item zero at 0
\item variable
\end{enumerate}
\end{enumerate}
\endgroup
\begingroup\small
\mxneed{18}\section*{MCQ 2: understand and choose}\label{C09-Q2}
Five single-answer questions for practice. Choose one answer per question; then explain one distractor. This chapter format does not reproduce the official EAM format.
\begin{enumerate}
\item $e^{2x}<e^3$ is equivalent to
\begin{enumerate}[label=\Alph*.,itemsep=2pt]
\item $x<3$
\item $x>\dfrac{3}{2}$
\item $x<\dfrac{3}{2}$
\item $x<\sqrt3$
\end{enumerate}
\item $u_n=e^{3n}$ has common ratio
\begin{enumerate}[label=\Alph*.,itemsep=2pt]
\item $e^3$
\item $3$
\item $e$
\item $3e$
\end{enumerate}
\item Tangent to $e^x$ at 0:
\begin{enumerate}[label=\Alph*.,itemsep=2pt]
\item $y=x$
\item $y=e x$
\item $y=1$
\item $y=x+1$
\end{enumerate}
\item Derivative of $xe^x$:
\begin{enumerate}[label=\Alph*.,itemsep=2pt]
\item $e^x$
\item $(x+1)e^x$
\item $xe^x$
\item $xe^{x-1}$
\end{enumerate}
\item $e^x\geq1+x$ has equality
\begin{enumerate}[label=\Alph*.,itemsep=2pt]
\item for all x
\item at 1 only
\item at 0 only
\item never
\end{enumerate}
\end{enumerate}
\subsection*{True or false? Justify.}\label{C09-VF}
\begin{enumerate}
\item $e^{2x}=2e^x$ for all x.
\item The function $t\mapsto e^{-3t}$ is decreasing.
\item A continuous rate k=\num{0,1} corresponds exactly to a \pc{10} increase in one unit of time.
\end{enumerate}
\begin{keypoint}[Self-assessment]I keep the rules for powers; I know the positivity; I differentiate with the factor; I separate model, exact calculation and simulation.

For each objective, indicate: achieved without help; achieved with help; to be reworked. Write one example and one mistake to avoid.\end{keypoint}
\endgroup
\ifmxteacher\mxinclude{corriges/C09}\fi
